\section*{Capítulo \ref{ch:g12:space} --- Vetores, retas e planos no espaço}

\begin{solution}{exo:g12:space:1}
$\vect{AB}(2, 1, -1)$ e $\vect{AC}(1, -1, 1)$ não são proporcionais
($\frac21 \neq \frac{1}{-1}$), logo os pontos não estão alinhados. Reta
$(AB)$:
\[
(x, y, z) = (1 + 2t,\; t,\; 2 - t), \qquad t \in \R .
\]
\end{solution}

\begin{solution}{exo:g12:space:2}
$2(x - 1) - (y - 1) + 3(z - 0) = 0$, \emph{isto é},
\[
2x - y + 3z - 1 = 0 .
\]
Para $B(0,2,1)$: $0 - 2 + 3 - 1 = 0$: sim, $B$ pertence ao plano.
\end{solution}

\begin{solution}{exo:g12:space:3}
$\vect{AB}(1,1,0)$, $\vect{AC}(1,0,1)$:
\[
\cos\widehat{A}
= \frac{\vect{AB}\cdot\vect{AC}}{\norm{\vect{AB}}\,\norm{\vect{AC}}}
= \frac{1}{\sqrt2\,\sqrt2} = \frac12,
\]
logo $\widehat A = 60^\circ$.
\end{solution}

\begin{solution}{exo:g12:space:4}
Pontos de $\mathcal D$: $(1 + t,\, 2 - t,\, 2t)$. Substituindo na \omterm{def:g10:algebra:equation}{equação}
de $\mathcal P$:
\[
2(1+t) + (2 - t) - 2t + 1 = 5 - t = 0 \iff t = 5 .
\]
Ponto de \omterm{def:g10:numbers:interunion}{interseção} único: $(6, -3, 10)$.
\end{solution}

\begin{solution}{exo:g12:space:5}
Os \omterm{def:g12:space:normal}{vetores normais} $(1, -1, 2)$ e $(2, 1, 1)$ não são \omterm{def:g12:space:collinear}{colineares}, de modo
que os planos se cortam segundo uma reta. Somando as duas \omterm{def:g10:algebra:equation}{equações}:
$3x + 3z = 5$, logo $x = \frac53 - z$; a primeira \omterm{def:g10:algebra:equation}{equação} dá então
$y = x + 2z - 1 = \frac23 + z$. Com $z = t$:
\[
(x, y, z) = \left(\frac53 - t,\; \frac23 + t,\; t\right), \qquad t \in \R,
\]
uma reta de \omterm{def:g11:vect:direction}{vetor diretor} $(-1, 1, 1)$. (As duas \omterm{def:g10:algebra:equation}{equações} se verificam
identicamente.)
\end{solution}

\begin{solution}{exo:g12:space:6}
As direções $\vec u_1(1, 2, -1)$ e $\vec u_2(1, 1, 2)$ não são \omterm{def:g12:space:collinear}{colineares},
logo as retas não são paralelas. Para que se cortassem, seria preciso
\[
1 + t = 2 + s, \qquad 2t = 1 + s, \qquad 3 - t = 1 + 2s .
\]
As duas primeiras dão $t = 1 + s$ e $2(1+s) = 1 + s$, logo $s = -1$,
$t = 0$; mas então a terceira diz $3 = -1$, absurdo. Sem ponto comum e não
paralelas: as retas são reversas.
\end{solution}

\begin{solution}{exo:g12:space:7}
\emph{1.} $\operatorname{dist} = \dfrac{\abs{2\cdot3 - 1 + 2\cdot1 - 3}}
{\sqrt{4 + 1 + 4}} = \dfrac{4}{3}$.

\emph{2.} $H = M_0 + t\,\vec n$ com $\vec n(2, -1, 2)$, escolhendo $t$ de
modo que $H \in \mathcal P$:
\[
2(3 + 2t) - (1 - t) + 2(1 + 2t) - 3 = 4 + 9t = 0 \iff t = -\frac49 .
\]
Logo $H\left(\frac{19}{9}, \frac{13}{9}, \frac19\right)$, e
\[
M_0H = \norm{t\,\vec n} = \frac49 \times 3 = \frac43,
\]
de acordo com a fórmula da distância.
\end{solution}

\begin{solution}{exo:g12:space:8}
\omterm{def:g10:coordgeom:system}{Coordenadas}: $G(1,1,1)$, e o plano $(BDE)$ passa por $B(1,0,0)$,
$D(0,1,0)$, $E(0,0,1)$.

\emph{1.} O plano $(BDE)$ tem \omterm{def:g10:algebra:equation}{equação} $x + y + z = 1$ (satisfeita pelos
três pontos, que não estão alinhados), de modo que $\vec n(1,1,1)$ é
normal a ele. Como $\vect{AG} = (1,1,1) = \vec n$, a diagonal $(AG)$ é
\omterm{def:g11:scal:orthogonal}{ortogonal} ao plano $(BDE)$.

\emph{2.} Distância de $A(0,0,0)$:
$\frac{\abs{0 + 0 + 0 - 1}}{\sqrt3} = \frac{1}{\sqrt3} = \frac{\sqrt3}{3}$.
A reta $(AG)$ é $(t, t, t)$; ela encontra o plano quando $3t = 1$: no
ponto $\left(\frac13, \frac13, \frac13\right)$, o baricentro do triângulo
$BDE$.
\end{solution}

\begin{solution}{exo:g12:space:9}
\emph{1.} $\vect{AB}(1, 0, -1)$, $\vect{AC}(-1, 1, 0)$,
$\vect{AD}(1, 1, 1)$. Se $\vect{AD} = a\vect{AB} + b\vect{AC}$, as
\omterm{def:g10:coordgeom:system}{coordenadas} dão $a - b = 1$, $b = 1$, $-a = 1$: as duas primeiras forçam
$a = 2$, contradizendo $a = -1$. Não \omterm{def:g12:space:collinear}{coplanares}.

\emph{2.} $\vect{BC}(-2, 1, 1)$, $\vect{BD}(0, 1, 2)$. Resolvendo
$\vec n \cdot \vect{BC} = \vec n \cdot \vect{BD} = 0$: de
$b + 2c = 0$, vem $b = -2c$; então $-2a - 2c + c = 0$ dá $c = -2a$. Com
$a = 1$: $\vec n(1, 4, -2)$. Plano por $B(2,1,0)$:
\[
x + 4y - 2z - 6 = 0 .
\]

\emph{3.} Distância de $A(1,1,1)$:
$\dfrac{\abs{1 + 4 - 2 - 6}}{\sqrt{21}} = \dfrac{3}{\sqrt{21}}$.
Área de $BCD$: $\norm{\vect{BC}}^2 = 6$, $\norm{\vect{BD}}^2 = 5$,
$\vect{BC}\cdot\vect{BD} = 3$, logo
\[
\text{Área} = \frac12\sqrt{6 \times 5 - 9} = \frac{\sqrt{21}}{2}.
\]

\emph{4.}
\[
V = \frac13 \times \frac{\sqrt{21}}{2} \times \frac{3}{\sqrt{21}}
= \frac12 .
\]
\end{solution}

\begin{solution}{pb:g12:space:1}
\textbf{1.} Plano: $x + y + z = 1$. Reta: $(t, t, t)$.
\omterm{def:g10:numbers:interunion}{Interseção}: $3t = 1$:
$\left(\frac13, \frac13, \frac13\right)$.

\textbf{2.} $\dfrac{\abs{0 + 0 + 0 - 1}}{\sqrt3} =
\dfrac{1}{\sqrt3} = \dfrac{\sqrt3}{3}$.

\textbf{3.} $1 - 1 + 0 = 0$: \omterm{def:g11:scal:orthogonal}{ortogonais}.

\textbf{4.} Mesma normal $(1,1,1)$, mas $\frac52 \neq 1$:
planos estritamente paralelos. Quanto à reta: a direção
$(1,1,0) \cdot (1,1,1) = 2 \neq 0$: não é paralela ao plano, de
modo que o atravessa em um único ponto.

\textbf{5.} \omterm{prop:g10:coordgeom:midpoint}{Pontos médios}: $\left(1, \frac12, 0\right)$,
$\left(\frac12, 1, 0\right)$, $\left(0, 1, \frac12\right)$,
$\left(0, \frac12, 1\right)$, $\left(\frac12, 0, 1\right)$,
$\left(1, 0, \frac12\right)$: cada um tem soma de \omterm{def:g10:coordgeom:system}{coordenadas}
$\frac32$.

\textbf{6.} A normal de $P$ é $(1,1,1)$, a direção de $(AG)$:
perpendicular. O centro
$\left(\frac12,\frac12,\frac12\right)$ tem soma de \omterm{def:g10:coordgeom:system}{coordenadas}
$\frac32$: está em $P$.

\textbf{7.} \omterm{prop:g10:coordgeom:midpoint}{Pontos médios} consecutivos diferem por \omterm{def:g11:vect:vector}{vetores} como
$\left(-\frac12, \frac12, 0\right)$: comprimento
$\frac{\sqrt2}{2}$ a cada vez --- seis lados iguais.

\textbf{8.} Em $\left(\frac12, 1, 0\right)$: vetores-aresta
$\left(\frac12, -\frac12, 0\right)$ e
$\left(-\frac12, 0, \frac12\right)$: \omterm{def:g12:space:dot}{produto escalar}
$-\frac14$, \omterm{def:g11:scal:dot}{normas} $\frac{\sqrt2}{2}$: $\cos = -\frac12$:
ângulo $120^\circ$. Lados iguais, ângulos iguais de
$120^\circ$: um hexágono regular.

\textbf{9.} Perímetro $6 \times \frac{\sqrt2}{2} = 3\sqrt2
\approx 4.24$. Área $6 \times \frac{\sqrt3}{4} \times
\frac12 = \frac{3\sqrt3}{4} \approx 1.30$: maior que uma face
(área $1$) --- uma fatia maior que qualquer lado da caixa; só
o retângulo diagonal ($\sqrt2 \approx 1.41$) a supera.

\textbf{10.} $(t,t,t)$ encontra $P$ em $t = \frac12$: o centro
--- que é o \omterm{prop:g10:coordgeom:midpoint}{ponto médio} de $[AG]$ ($A$ em $t = 0$, $G$ em
$t = 1$).

\textbf{11.} Para $c = \frac12$: o plano corta as três arestas
do canto de $A$, em $\left(\frac12,0,0\right)$,
$\left(0,\frac12,0\right)$, $\left(0,0,\frac12\right)$: um
triângulo equilátero de lado $\frac{\sqrt2}{2}$. À medida que
$c$ cresce, o triângulo cresce, seus cantos vão sendo
truncados até virar um hexágono (regular exatamente em
$c = \frac32$), e depois a figura encolhe simetricamente até
um triângulo no canto de $G$: a metamorfose da lapidação.

\textbf{12.} Cada lado da seção é a \omterm{def:g10:numbers:interunion}{interseção} do plano de
corte com uma \emph{face} do cubo, e um plano encontra uma
face (um polígono plano) em no máximo um segmento: no máximo
$6$ lados. Sete é impossível; de três a seis, todos ocorrem
(as questões 8 e 11 mostram dois casos).

\textbf{13.} Base $ABD$: triângulo retângulo de área
$\frac12$; altura $AE = 1$: volume
$\frac13 \times \frac12 \times 1 = \frac16$.

\textbf{14.} Cada par entre $B, D, E, G$ difere em exatamente
duas \omterm{def:g10:coordgeom:system}{coordenadas} por $1$: as seis distâncias valem $\sqrt2$
--- regular. O cubo se decompõe em $BDEG$ mais quatro
tetraedros de canto congruentes a $ABDE$: volume
$1 - 4 \times \frac16 = \frac13$.

\textbf{15.} Plano $(BDE)$: $x + y + z = 1$; centro
$\left(\frac12,\frac12,\frac12\right)$: distância
$\frac{\abs{\frac32 - 1}}{\sqrt3} = \frac{1}{2\sqrt3} =
\frac{\sqrt3}{6}$.

\textbf{16.} $\vect{AG} = (1,1,1)$,
$\vect{BH} = (-1,1,1)$: $\cos\theta = \frac{-1 + 1 + 1}{3} =
\frac13$: $\theta \approx 70.5^\circ$, suplemento
$\approx 109.5^\circ$. Quatro pares de elétrons se repelem e
se afastam o \omterm{def:g10:functions:extrema}{máximo} possível em torno do carbono: o ótimo são
os cantos do tetraedro regular, cujas direções do centro aos
vértices se encontram exatamente nesse ângulo --- o diamante é
a questão 14 cristalizada.

\textbf{17.} $(1, t, t)$ em $P$: $1 + 2t = \frac32$:
$t = \frac14$: o ponto $\left(1, \frac14, \frac14\right)$.

\textbf{18.} $(AB)$: pontos $(t, 0, 0)$; $(EG)$: pontos
$(s, s, 1)$. Encontrarem-se exigiria $0 = 1$ na terceira
coordenada: nunca. Serem paralelas exigiria $(1,0,0)$ e
$(1,1,0)$ \omterm{def:g12:space:collinear}{colineares}: não. Nem se encontram nem são paralelas:
reversas --- dois corredores em andares diferentes.

\textbf{19.} $B + t(1,1,1) = (1 + t, t, t)$ em $P$:
$1 + 3t = \frac32$: $t = \frac16$: projeção
$\left(\frac76, \frac16, \frac16\right)$, cuja soma de
\omterm{def:g10:coordgeom:system}{coordenadas} é $\frac32$: está em $P$, como se pedia.

\textbf{20.} As retas paramétricas atravessam (questões 1, 10,
17); os \omterm{def:g12:space:normal}{vetores normais} medem ângulos e distâncias (questões
2, 8, 15, 16); as \omterm{def:g10:algebra:equation}{equações} de planos recortam seções (a
metamorfose do triângulo ao hexágono). E o pátio de recreio
retribuiu a visita: um hexágono regular numa caixa de
quadrados, um tetraedro regular nos cantos, o ângulo do
diamante e duas retas que se ignoram --- geometria que só o
espaço oferece.
\end{solution}
